\documentclass[a4paper, 12pt]{article}

\usepackage{utils}

\renewcommand*{\today}{10 February 2026}

\begin{document}

\hotbox{Algèbre bilinéaire}{CM 1}{\today}

\section{Formes bilinéaires et formes quadratiques}

Dans cette section, nous ne considérons par $\K$ que les corps commutatifs de caractéristique différente de $2$.

Soit $E$ et $F$ des $\K$-espaces vectoriels. On notera $\mathcal{L}(E, F)$ l'espace des applications linéaires de $E$ dans $F$.
Si $E$ est un $\K$-espace vectoriel, on notera $E^* = \mathcal{L}(E, \K)$ son dual.

\begin{definition}
    Soit $E$, $F$ et $G$ des $\K$-espaces vectoriels et $\varphi: E \times F \to G$ une application.

    On dit que l'application $\varphi$ est une \textbf{application bilinéaire} lorsque, pour tout $u, u' \in E$, $v, v' \in F$ et $\lambda \in \K$, on a
    \begin{align*}  
        \varphi(u + \lambda u', v) = \varphi(u, v) + \lambda \varphi(u', v) \\
        \varphi(u, v + \lambda v') = \varphi(u, v) + \lambda \varphi(u, v') \\
    \end{align*}
\end{definition}

\begin{proposition}{}{bililiniearite}
    L'application $\varphi: E \times F \to G$ est bilinéaire si et seulement si,
    $$
    \forall v \in F, \quad \varphi(\cdot, v) \in \mathcal{L}(E, G) \quad \text{et} \quad \forall u \in E, \quad \varphi(u, \cdot) \in \mathcal{L}(F, G)
    $$
    où pour tout $v \in F$,
    $$
    \funcdef{\varphi(\cdot, v)}{E}{G}{u}{\varphi(u, v)}
    $$
    et pour tout $u \in E$,
    $$
    \funcdef{\varphi(u, \cdot)}{F}{G}{v}{\varphi(u, v)}
    $$
\end{proposition}

\begin{demonstration}
    \begin{description}
        \item[Sens direct.]
        On suppose que $\varphi$ est bilinéaire. Soit $u \in E$. Alors, pour tout $v, v' \in F$ et $\lambda \in \K$, on a
        \begin{align*}
            \varphi(u, \cdot)(v + \lambda v') &= \varphi(u, v + \lambda v') \\
            &\underset{\text{bilinéarité}}{=} \varphi(u, v) + \lambda \varphi(u, v') \\
            &= \varphi(u, \cdot)(v) + \lambda \varphi(u, \cdot)(v')
        \end{align*}
        Donc $\varphi(u, \cdot)$ est linéaire. Même raisonnement pour $\varphi(\cdot, v)$.

        \item[Sens réciproque.]
        On suppose que pour tout $u \in E$, $v \in F$, $\varphi(\cdot, v) \in \mathcal{L}(E, G)$ et $\varphi(u, \cdot) \in \mathcal{L}(F, G)$. Soient $u, u' \in E$, $v, v' \in F$ et $\lambda \in \K$. Alors
        \begin{align*}
            \varphi(u + \lambda u', v) &= \varphi(\cdot, v)(u + \lambda u') \\
            &= \varphi(\cdot, v)(u) + \lambda \varphi(\cdot, v)(u') \\
            &= \varphi(u, v) + \lambda \varphi(u', v)
        \end{align*}
        et
        \begin{align*}
            \varphi(u, v + \lambda v') &= \varphi(u, \cdot)(v + \lambda v') \\
            &= \varphi(u, \cdot)(v) + \lambda \varphi(u, \cdot)(v') \\
            &= \varphi(u, v) + \lambda \varphi(u, v')
        \end{align*}
        Donc $\varphi$ est bilinéaire.
    \end{description}
\end{demonstration}

\begin{definition}
    Soit $\varphi: E \times F \to \K$ une application bilinéaire.
    \begin{itemize}[label=--]
        \item On dit que $\varphi$ est une \textbf{forme bilinéaire} sur $E \times F$.
        \item Lorsque, de plus $E = F$ on dit que $\varphi$ est une \textbf{forme bilinéaire} sur $E$.
    \end{itemize}
\end{definition}

\begin{proposition}{}{gliniearite}
    Soient $E$ et $F$ deux $\K$-espaces vectoriels et soit $\varphi: E \times F \to \K$ une forme bilinéaire. Alors, les applications
    $$
    \funcdef{g_\varphi}{E}{F^*}{u}{\varphi(u, \cdot)} \quad \text{et} \quad \funcdef{d_\varphi}{F}{E^*}{v}{\varphi(\cdot, v)}
    $$
    sont linéaires.
\end{proposition}

\begin{demonstration}
    On suppose que $\varphi$ est une forme bilinéaire sur $E \times F$. Soient $u, u' \in E$, $v, v' \in F$ et $\lambda \in \K$. Alors
    \begin{align*}
        g_\varphi(u + \lambda u') &= \varphi(u + \lambda u', \cdot) \\
        &= \varphi(u, \cdot) + \lambda \varphi(u', \cdot) \quad \text{(*)} \\
        &= g_\varphi(u) + \lambda g_\varphi(u')
    \end{align*}
    (*) est vérifié par la proposition \ref{propal:bililiniearite}.

    Donc $g_\varphi$ est linéaire. Même raisonnement pour $d_\varphi$.
\end{demonstration}

\subsection{Espace des formes bilinéaires}

Dans cette section, on considère que $E$ et $F$ sont des $\K$-espaces vectoriels de dimension finie $p$ et $n$ respectivement et leurs bases $\mathscr{B} = (e_1, \dots, e_p)$ et $\mathscr{B'} = (e_1', \dots, e_n')$ respectivement.
\vspace{1cm}

\begin{proposition}{}{}
    L'ensemble des formes bilinéaires sur $E \times F$ est un $\K$-espace vectoriel, noté $\mathcal{L}_2(E, F)$,
    isomorphe à $\mathcal{L}(E, F^*)$ et à $\mathcal{L}(F, E^*)$.
\end{proposition}

\begin{demonstration}
    On définie l'application $\funcdef{\Phi}{\mathcal{L}_2(E, F)}{\mathcal{L}(E, F^*)}{\varphi}{g_\varphi}$.

    \begin{enumerate}
        \item $\Phi$ est bilinéaire, d'après la proposition \ref{propal:bililiniearite} $\varphi(u, \cdot)$ est linéaire pour tout $u \in E$.
        \item $\Phi$ est bilinéaire, d'après la proposition \ref{propal:gliniearite} $g_\varphi$ est linéaire pour tout $\varphi \in \mathcal{L}_2(E, F)$.
        
        \item
        \begin{align*}
            \Phi(\varphi) = 0 &\implies \forall u \in E, \quad g_\varphi(u) = 0 \\
            &\implies \forall (u, v) \in E \times F, \quad \varphi(u, v) = 0 \\
            &\implies \varphi = 0
        \end{align*}
        Donc $\Phi$ est injective.
        \begin{hotwarn}
            Faire la surjectivité
        \end{hotwarn}
    \end{enumerate}
    On a donc que $\Phi$ est un isomorphisme de $\K$-espaces vectoriels. De même, on peut construire un isomorphisme de $\mathcal{L}_2(E, F)$ vers $\mathcal{L}(F, E^*)$.
\end{demonstration}

\subsection{Écriture matricielle des formes bilinéaires}

\begin{remarque}
    On suppose que $E$ est de dimension finie $p \in \N^*$ et $F$ est de dimension finie $n \in \N^*$.

    Comme $dim \mathcal{L}(E, F^*) = pn$, on a $dim \mathcal{L}_2(E, F) = pn$.
\end{remarque}

\begin{definition}
    Soit $\mathscr{B} = (e_1, \dots, e_p)$ une base de $E$ et $\mathscr{B'} = (e_1', \dots, e_n')$ une base de $F$.
    Soit $\varphi \in \mathcal{L}_2(E, F)$. Soit $M = (\varphi(e_i, e_j'))_{1 \leq i \leq p, 1 \leq j \leq n}$.
    On dit que cette matrice est la matrice de $\varphi$ dans les bases $\mathscr{B}$ et $\mathscr{B'}$.
    On note alors $M = Mat_{\varphi; \mathscr{B}, \mathscr{B'}} \in \mathcal{M}_\K(p, n)$
    $$
    \begin{pmatrix}
        \varphi(e_1, e_1') & \varphi(e_1, e_2') & \dots & \varphi(e_1, e_n') \\
        \varphi(e_2, e_1') & \varphi(e_2, e_2') & \dots & \varphi(e_2, e_n') \\
        \vdots & \vdots & \ddots & \vdots \\
        \varphi(e_p, e_1') & \varphi(e_p, e_2') & \dots & \varphi(e_p, e_n')
    \end{pmatrix}
    $$
\end{definition}

\begin{theoreme}{Représentation matricielle}{}
    Soient $E$ et $F$ deux $\K$-espaces vectoriels de dimension finie $p$ et $n$ respectivement.
    Soient $\mathscr{B} = (e_1, \dots, e_p)$ une base de $E$ et $\mathscr{B'} = (e_1', \dots, e_n')$ une base de $F$.

    Pour toute forme bilinéaire $\varphi \in \mathcal{L}_2(E, F)$, il existe une \textbf{unique} matrice $M \in \mathcal{M}_\K(p, n)$ telle que:
    \begin{itemize}
        \item pour tout $i \in \llbracket 1, p \rrbracket$ et $j \in \llbracket 1, n \rrbracket$, on a $m_{i, j} = \varphi(e_i, e_j')$.
        \item pour tout $u \in E$ et $v \in F$, on a
        
        $$\varphi(u, v) = X^T M Y = Y^T X^T M$$

        où $X$ et $Y$ sont les vecteurs colonnes des coordonnées de $u$ dans $\mathscr{B}$ et de $v$ dans $\mathscr{B'}$ respectivement.
    \end{itemize}
\end{theoreme}

\begin{demonstration}
    Soit $\varphi \in \mathcal{L}_2(E, F)$. Soient $\mathscr{B} = (e_1, \dots, e_p)$ une base de $E$ et $\mathscr{B'} = (e_1', \dots, e_n')$ une base de $F$.
    \begin{description}
        \item[Existence.] Soient $u \in E$ et $v \in F$. On décompose $u$ et $v$ dans leurs bases respectives, alors
        $$
        u = \sum_{i=1}^p x_i e_i \quad \text{et} \quad v = \sum_{j=1}^n y_j e_j'
        $$
        Par la bilinéarité de $\varphi$, on a
        \begin{align*}
            \varphi(u, v) &= \varphi(\sum_{i=1}^p x_i e_i, \sum_{j=1}^n y_j e_j') \\
            &= \sum_{i=1}^p \sum_{j=1}^n x_i y_j \varphi(e_i, e_j')
        \end{align*}
        En posant $m_{i, j} = \varphi(e_i, e_j')$ pour tout $i \in \llbracket 1, p \rrbracket$ et $j \in \llbracket 1, n \rrbracket$, on a
        $$
        \varphi(u, v) = \sum_{i=1}^p x_i \left( \sum_{j=1}^n m_{i, j} y_j \right) = \sum_{i=1}^p x_i (M Y)_i = X^T M Y
        $$

        \item[Unicité.] Supposons qu'il existe une matrice $B \in \mathcal{M}_\K(p, n)$ telle que pour tout $u \in E$ et $v \in F$, on a $\varphi(u, v) = X^T B Y$.
        Ainsi, pour tout $i \in \llbracket 1, p \rrbracket$ et $j \in \llbracket 1, n \rrbracket$, on a
        \begin{align*}
            m_{i, j} = \varphi(e_i, e_j') &= E_i^T M E_j' \\
            &= E_i^T B E_j' \\
            &= b_{i, j}
        \end{align*}
        où $E_i$ et $E_j'$ sont les vecteurs colonnes de coordonnées de $e_i$ dans $\mathscr{B}$ et de $e_j'$ dans $\mathscr{B'}$ respectivement. Donc $M = B$.
    \end{description}
\end{demonstration}

\begin{proposition}{}{}
    L'application de $\mathcal{L}_2(E, F)$ sur $\mathcal{M}_\K(p, n)$ définie par $\varphi \mapsto Mat_{\varphi; \mathscr{B}, \mathscr{B'}}$ est un isomorphisme de $\K$-espaces vectoriels.
\end{proposition}

\begin{theoreme}{Formule de changement de base}{}
    Soit $\varphi$ une forme bilinéaire sur $E \times F$.
    On considère:
    \begin{itemize}
        \item $\mathscr{B}$ et $\mathscr{B'}$ les bases "anciennes" de $E$ et $F$ respectivement
        \item $\mathscr{B}_1$ et $\mathscr{B'}_1$ les bases "nouvelles" de $E$ et $F$ respectivement
    \end{itemize}
    Soient $P$ la matrice de passage de $\mathscr{B}$ vers $\mathscr{B}_1$ et $Q$ la matrice de passage de $\mathscr{B'}$ vers $\mathscr{B'}_1$. Alors
    $$
    Mat_{\varphi; \mathscr{B}_1, \mathscr{B'}_1} = P^T Mat_{\varphi; \mathscr{B}, \mathscr{B'}} Q
    $$
\end{theoreme}

\begin{demonstration}
    Par définition de $P$ et $Q$, on a $X = P X_1$ et $Y = Q Y_1$.

    Par définition de $Mat_{\varphi; \mathscr{B}_1, \mathscr{B'}_1}$, pour tout $u \in E$ et $v \in F$, on a
    $$
    \varphi(u, v) = X_1^T Mat_{\varphi; \mathscr{B}_1, \mathscr{B'}_1} Y_1
    $$
    De plus, par définition de $Mat_{\varphi; \mathscr{B}, \mathscr{B'}}$, pour tout $u \in E$ et $v \in F$, on a
    \begin{align*}
    \varphi(u, v) &= X^T Mat_{\varphi; \mathscr{B}, \mathscr{B'}} Y \\
    &= (P X_1)^T Mat_{\varphi; \mathscr{B}, \mathscr{B'}} (Q Y_1) \\
    &= X_1^T \left( P^T Mat_{\varphi; \mathscr{B}, \mathscr{B'}} Q \right) Y_1
    \end{align*}
    Donc par unicité de la représentation matricielle, on a $Mat_{\varphi; \mathscr{B}_1, \mathscr{B'}_1} = P^T Mat_{\varphi; \mathscr{B}, \mathscr{B'}} Q$.
\end{demonstration}

\subsection{Rang d'une forme bilinéaire}

\begin{definition}
    On suppose que $E$ et $F$ sont de dimension finie. Soit $\varphi$ une forme bilinéaire sur $E \times F$. 
    Alors les applications linéaires $g_\varphi: E \to F^*$ et $d_\varphi: F \to E^*$ sont de même rang, appelé le \textbf{rang} de $\varphi$ et noté $rg(\varphi)$.

    De plus, le rang de $\varphi$ est égal au rang de la matrice de $\varphi$ dans n'importe quelle base de $E$ et n'importe quelle base de $F$.
\end{definition}

\subsection{Formes bilinéaires non dégénérées}

\begin{definition}
    Une forme bilinéaire $\varphi$ sur $E \times F$ est dite \textbf{non dégénérée} lorsque les applications linéaires $g_\varphi$ et $d_\varphi$ sont injectives. Une forme bilinéaire $\varphi$ sur $E \times F$ qui n'est pas dégénérée est dit \textbf{dégénérée}.
\end{definition}

\begin{proposition}{}{}
    On suppose que les deux espaces vectoriels $E$ et $F$ sont de dimension finie.
    Soit $\varphi \in \mathcal{L}_2(E, F)$ et $M$ la matrice de $\varphi$ dans une base de $E$ et une base de $F$. Alors, les propriétés suivantes sont équivalentes:
    \begin{enumerate}
        \item $\varphi$ est non dégénérée
        \item $g_\varphi$ est bijective
        \item $d_\varphi$ est bijective
        \item $\forall u \in E, \left( \forall v \in F, \quad \varphi(u, v) = 0 \implies u = 0 \right)$ et $dim E = dim F$
        \item $\forall v \in F, \left( \forall u \in E, \quad \varphi(u, v) = 0 \implies v = 0 \right)$ et $dim E = dim F$
        \item $rang(\varphi) = dim E = dim F$
        \item $dim E = dim F$ et $det(M) \neq 0$
    \end{enumerate}
\end{proposition}

\subsection{Applications bilinéaires symétriques et antisymétriques}

\begin{definition}
    Soit $E$ et $G$ des $\K$-espaces vectoriels et soit $\varphi: E \times E \to G$.
    \begin{itemize}[label=--]
        \item $\varphi$ est dite \textbf{symétrique} lorsque, pour tout $(u, v) \in E \times E$, on a $\varphi(u, v) = \varphi(v, u)$.
        \item $\varphi$ est dite \textbf{antisymétrique} lorsque, pour tout $(u, v) \in E \times E$, on a $\varphi(u, v) = -\varphi(v, u)$.
        \item $\varphi$ est dite \textbf{alternée} lorsque, pour tout $u \in E$, on a $\varphi(u, u) = 0$.
    \end{itemize}
\end{definition}

\begin{proposition}{}{}
    Si de plus, $\varphi$ est bilinéaire, alors $\varphi$ est alternée si et seulement si $\varphi$ est antisymétrique.
    (car la caractéristique de $\K$ est différente de $2$)
\end{proposition}

\begin{demonstration}
    \begin{description}
        \item[Sens direct.]
        Supposons que $\varphi$ est antisymétrique. Soient $x, y \in E$. Alors
        \begin{hotwarn}
            À faire
        \end{hotwarn}

        \item[Sens réciproque.]
        Supposons que $\varphi$ est alternée. Soient $x, y \in E$. Alors
        \begin{align*}
            &\varphi(x + y, x + y) = 0 \\
            \implies &\varphi(x, x) + \varphi(y, y) + 2\varphi(x, y) = 0 \\
            \implies &\varphi(x, y) = -\varphi(y, x)
        \end{align*}
        Donc $\varphi$ est antisymétrique.
    \end{description}
\end{demonstration}

\begin{proposition}{}{}
    L'ensemble des formes bilinéaires symétriques (resp. antisymétriques) sur $E$ est un espace vectoriel, noté $\mathcal{S}_2(E)$ (resp. $\mathcal{A}_2(E)$), isomorphe à un sous-espace de $\mathcal{L}(E, E^*)$.
    donc
    $$
    \mathcal{L}_2(E) = \mathscr{S}_2(E) \oplus \mathscr{A}_2(E)
    $$
\end{proposition}

\begin{demonstration}
    Soit $\varphi \in \mathcal{L}_2(E)$.

    $$
    \varphi(x, y) = \frac{1}{2}(\varphi(x, y) + \varphi(y, x)) + \frac{1}{2}(\varphi(x, y) - \varphi(y, x))
    $$

    Soit $\varphi_e \in \mathscr{S}_2(E) \cap \mathscr{A}_2(E)$. Alors, pour tout $x, y \in E$, on a
    \begin{align*}
        &\varphi_e(x, y) = \varphi_e(y, x) \\
        &\varphi_e(x, y) = -\varphi_e(y, x)
    \end{align*}
    Donc $\varphi_e(x, y) = 0$ pour tout $x, y \in E$ et $\varphi_e$ est la forme bilinéaire nulle.
\end{demonstration}

\end{document}